% ============================================================
\chapter{Sobolev Spaces}
\label{ch:sobolev}
% ============================================================

Classical partial differential equations demand that solutions be smooth enough for derivatives to make pointwise sense. But many physical problems --- vibrating membranes with corners, turbulent flows, composite materials --- produce solutions that are not differentiable in the classical sense. Sergei Sobolev, in the 1930s, resolved this dilemma by introducing \emph{weak derivatives}: a notion of differentiation that does not require pointwise regularity but only that the ``integration by parts formula'' hold in the distributional sense. The \emph{Sobolev spaces} $W^{k,p}(\Omega)$, equipped with natural norms, are the function spaces where weak solutions of PDEs live. Today, all of modern PDE theory --- variational methods, finite elements, regularity theory --- rests on this framework.

\section{Weak derivatives}

\begin{definition}[Weak derivative]
\index{weak derivative}
Let $\Omega \subset \R^n$ be an open set and $u \in L^1_{\mathrm{loc}}(\Omega)$.
We say that $v \in L^1_{\mathrm{loc}}(\Omega)$ is the \emph{weak derivative}
$\partial^\alpha u$ (for a multi-index $\alpha$) if:
\[
    \int_\Omega v\, \varphi\, dx = (-1)^{\abs{\alpha}} \int_\Omega u\, \partial^\alpha \varphi\, dx
    \quad \forall\, \varphi \in C_c^\infty(\Omega).
\]
\end{definition}

\begin{intuition}
The weak derivative generalizes the classical derivative to functions that are not
necessarily differentiable in the usual sense. The formula is based on integration
by parts: if $u$ is $C^k$, the weak derivative coincides with the classical one.
The key point is that a weak derivative can exist even if $u$ has singularities.
\end{intuition}

\begin{example}[Absolute value]
The function $u(x) = \abs{x}$ on $\R$ has weak derivative
$u'(x) = \operatorname{sgn}(x)$ (the sign function). Indeed, for all
$\varphi \in C_c^\infty(\R)$:
\[
    \int_\R \abs{x} \varphi'(x)\, dx = -\int_\R \operatorname{sgn}(x)\, \varphi(x)\, dx.
\]
\end{example}

\begin{example}[Heaviside function]
The Heaviside function $H(x) = \ind_{(0,\infty)}(x)$ does not have a weak
derivative in $L^1_{\mathrm{loc}}(\R)$: its distributional derivative is the
Dirac mass $\delta_0$, which is not a function.
\end{example}

\section{Sobolev spaces $W^{k,p}$}

\begin{definition}[Sobolev space $W^{k,p}(\Omega)$]
\index{Sobolev space}
Let $\Omega \subset \R^n$ be open, $k \in \N$ and $1 \leq p \leq \infty$.
The \emph{Sobolev space} $W^{k,p}(\Omega)$ is:
\[
    W^{k,p}(\Omega) = \{ u \in L^p(\Omega) : \partial^\alpha u \in L^p(\Omega),\;
    \forall\, \abs{\alpha} \leq k \},
\]
equipped with the norm:
\[
    \norm{u}_{W^{k,p}} = \left( \sum_{\abs{\alpha} \leq k}
    \norm{\partial^\alpha u}_{L^p}^p \right)^{1/p}
    \quad (1 \leq p < \infty).
\]
For $p = 2$, we write $H^k(\Omega) = W^{k,2}(\Omega)$.
\end{definition}

\begin{theorem}[Completeness]
The space $W^{k,p}(\Omega)$ is a Banach space. For $p = 2$, $H^k(\Omega)$ is a
Hilbert space with inner product:
\[
    \inner{u}{v}_{H^k} = \sum_{\abs{\alpha} \leq k} \int_\Omega
    \partial^\alpha u \cdot \overline{\partial^\alpha v}\, dx.
\]
\end{theorem}

\begin{definition}[Space $W_0^{k,p}(\Omega)$]
\index{$W_0^{k,p}$ space}
The space $W_0^{k,p}(\Omega)$ is the closure of $C_c^\infty(\Omega)$ in
$W^{k,p}(\Omega)$. Functions in $W_0^{k,p}$ are those that ``vanish on the
boundary'' in the sense of traces. We write $H_0^k(\Omega) = W_0^{k,2}(\Omega)$.
\end{definition}

\section{Density of smooth functions}

\begin{theorem}[Meyers-Serrin]
\index{Meyers-Serrin theorem}
The space $C^\infty(\Omega) \cap W^{k,p}(\Omega)$ is dense in $W^{k,p}(\Omega)$
for $1 \leq p < \infty$.
\end{theorem}

\begin{proof}[Proof sketch]
One uses a partition of unity $\{\rho_j\}$ and convolution regularizations
$(\rho_j u) * \eta_{\varepsilon_j}$ with mollifiers $\eta_\varepsilon$. By
choosing $\varepsilon_j$ small enough for each $j$, one approximates $u$ in
$W^{k,p}$.
\end{proof}

\begin{theorem}[Density of $C^\infty(\overline{\Omega})$]
If $\Omega$ is a bounded open set with $C^1$ boundary, then
$C^\infty(\overline{\Omega})$ is dense in $W^{k,p}(\Omega)$ for $1 \leq p < \infty$.
\end{theorem}

\section{Sobolev embedding theorems}

\begin{theorem}[Sobolev embedding]
\label{thm:sobolev_embedding}
\index{Sobolev embedding}
Let $\Omega \subset \R^n$ be a bounded open set with Lipschitz boundary,
$k \in \N^*$ and $1 \leq p < \infty$.
\begin{enumerate}
    \item If $kp < n$: $W^{k,p}(\Omega) \hookrightarrow L^{p^*}(\Omega)$
          with $p^* = \frac{np}{n - kp}$ (Sobolev exponent).
    \item If $kp = n$: $W^{k,p}(\Omega) \hookrightarrow L^q(\Omega)$
          for all $q \in [p, \infty)$.
    \item If $kp > n$: $W^{k,p}(\Omega) \hookrightarrow C^{0,\gamma}(\overline{\Omega})$
          with $\gamma = k - n/p$ if $k - n/p < 1$.
\end{enumerate}
\end{theorem}

\begin{warning}[title=\textbf{The critical exponent}]
The embedding $W^{1,p} \hookrightarrow L^{p^*}$ is \emph{continuous} but
\emph{not compact}. The Rellich-Kondrachov theorem gives compactness for
exponents strictly less than $p^*$.
\end{warning}

\begin{theorem}[Rellich-Kondrachov]
\label{thm:rellich}
\index{Rellich-Kondrachov theorem}
Let $\Omega$ be a bounded open set with Lipschitz boundary.
\begin{enumerate}
    \item If $kp < n$: $W^{k,p}(\Omega) \hookrightarrow\hookrightarrow L^q(\Omega)$
          for all $q < p^*$ (compact embedding).
    \item If $kp > n$: $W^{k,p}(\Omega) \hookrightarrow\hookrightarrow C(\overline{\Omega})$
          (compact embedding).
\end{enumerate}
\end{theorem}

\begin{example}[Case $n = 1$]
In dimension $1$, for $p = 2$ and $k = 1$: $kp = 2 > n = 1$, so
$H^1(\Omega) \hookrightarrow C(\overline{\Omega})$. Every $H^1$ function in
dimension $1$ is continuous (after modification on a null set).
\end{example}

\begin{example}[Case $n = 3$, $k = 1$, $p = 2$]
$kp = 2 < n = 3$, so $p^* = \frac{3 \cdot 2}{3 - 2} = 6$.
We have $H^1(\Omega) \hookrightarrow L^6(\Omega)$ (continuous) and
$H^1(\Omega) \hookrightarrow\hookrightarrow L^q(\Omega)$ for $q < 6$ (compact).
\end{example}

\section{Trace theorem}

\begin{theorem}[Trace]
\label{thm:trace}
\index{trace theorem}
Let $\Omega$ be a bounded open set with $C^1$ boundary. The restriction operator
$\gamma_0 : C^\infty(\overline{\Omega}) \to C^\infty(\partial\Omega)$,
$\gamma_0(u) = u|_{\partial\Omega}$, extends to a continuous surjective linear
operator:
\[
    \gamma_0 : W^{1,p}(\Omega) \to W^{1-1/p, p}(\partial\Omega),
    \quad 1 < p < \infty.
\]
Moreover, $\ker \gamma_0 = W_0^{1,p}(\Omega)$.
\end{theorem}

\begin{remark}
This theorem justifies the Dirichlet condition $u = g$ on $\partial\Omega$ in the
Sobolev framework: one requires $\gamma_0(u) = g$ in the trace sense, even if $u$
is not continuous.
\end{remark}

\section{Poincar\'e inequality}

\begin{theorem}[Poincar\'e inequality]
\label{thm:poincare}
\index{Poincar\'e inequality}
Let $\Omega \subset \R^n$ be a bounded open set and $1 \leq p < \infty$. There
exists a constant $C = C(\Omega, p, n) > 0$ such that:
\[
    \norm{u}_{L^p(\Omega)} \leq C \norm{\nabla u}_{L^p(\Omega)}
    \quad \forall\, u \in W_0^{1,p}(\Omega).
\]
\end{theorem}

\begin{proof}
By contradiction. Suppose there exists a sequence $(u_k)$ in $W_0^{1,p}(\Omega)$
with $\norm{u_k}_{L^p} = 1$ and $\norm{\nabla u_k}_{L^p} \to 0$. By
Rellich-Kondrachov, $(u_k)$ has a subsequence converging in $L^p$. The limit $u$
satisfies $\nabla u = 0$ a.e., so $u$ is constant. Since $u \in W_0^{1,p}$, we
have $u = 0$, contradicting $\norm{u}_{L^p} = 1$.
\end{proof}

\begin{corollary}
On $W_0^{1,p}(\Omega)$ with $\Omega$ bounded, the seminorm $\norm{\nabla u}_{L^p}$
is equivalent to the full norm $\norm{u}_{W^{1,p}}$.
\end{corollary}

\begin{theorem}[Poincar\'e-Wirtinger inequality]
\index{Poincar\'e-Wirtinger inequality}
Let $\Omega$ be a bounded connected open set with Lipschitz boundary. For
$1 \leq p < \infty$:
\[
    \norm{u - \bar{u}}_{L^p(\Omega)} \leq C \norm{\nabla u}_{L^p(\Omega)}
    \quad \forall\, u \in W^{1,p}(\Omega),
\]
where $\bar{u} = \frac{1}{\abs{\Omega}} \int_\Omega u\, dx$ is the mean of $u$.
\end{theorem}

\begin{keyformulas}[title=\textbf{Summary of Sobolev embeddings ($k=1$)}]
\begin{center}
\renewcommand{\arraystretch}{1.3}
\begin{tabular}{lcl}
\textbf{Condition} & \textbf{Embedding} & \textbf{Type} \\
\hline
$p < n$ & $W^{1,p} \hookrightarrow L^{np/(n-p)}$ & continuous \\
$p < n$ & $W^{1,p} \hookrightarrow\hookrightarrow L^{q}$, $q < \frac{np}{n-p}$ & compact \\
$p = n$ & $W^{1,p} \hookrightarrow L^q$, $\forall q < \infty$ & continuous \\
$p > n$ & $W^{1,p} \hookrightarrow C^{0,1-n/p}$ & continuous \\
$p > n$ & $W^{1,p} \hookrightarrow\hookrightarrow C(\overline{\Omega})$ & compact \\
\end{tabular}
\end{center}
\end{keyformulas}

\section{Exercises}

\begin{exercise}
Show that $u(x) = \abs{x}^{-\alpha}$ belongs to $W^{1,p}(B(0,1))$ (unit ball in
$\R^n$) if and only if $(\alpha + 1)p < n$.
\end{exercise}

\begin{exercise}
Let $\Omega = (0,1)$ and $u(x) = x^\beta$ for $\beta > 0$. Determine for which
values of $\beta$ and $k$ we have $u \in H^k(\Omega)$.
\end{exercise}

\begin{exercise}
Prove the Sobolev inequality in dimension $1$: for $u \in W^{1,1}(\R)$,
$\norm{u}_{L^\infty} \leq \frac{1}{2}\norm{u'}_{L^1}$.
\end{exercise}

\begin{exercise}
Show that $H^1(\R^n)$ does not embed into $L^\infty(\R^n)$ for $n \geq 2$.
\emph{Hint:} construct a logarithmic counterexample.
\end{exercise}

\begin{exercise}
Let $\Omega$ be a bounded open subset of $\R^n$ with smooth boundary. Show that
the Poincar\'e constant $C_P$ satisfies $C_P \leq \frac{\operatorname{diam}(\Omega)}{2}$
when $p = 2$.
\end{exercise}

\begin{exercise}
Show that if $u \in H^1_0(\Omega)$ and $f \in L^2(\Omega)$ satisfy $-\Delta u = f$
in the weak sense, then $\norm{u}_{H^1_0} \leq C\norm{f}_{L^2}$ with a constant
depending only on $\Omega$.
\end{exercise}
